El campo de la Interpretación y Compilación de Lenguajes de Programación constituye una pieza fundamental en la formación de ingenieros de sistemas; sin embargo, sus contenidos suelen ser percibidos como abstractos y complejos por los estudiantes.

La principal dificultad radica en el alto nivel de abstracción de conceptos esenciales, tales como el análisis sintáctico, la gestión de Árboles de Sintaxis Abstracta (AST) y la representación del ambiente de ejecución de programas. Esta complejidad ha motivado la búsqueda de enfoques pedagógicos alternativos. Por ejemplo, la literatura especializada ha resaltado la efectividad del aprendizaje por proyectos para mejorar la comprensión de conceptos \cite{liu2017evaluation}, así como la utilidad de lenguajes simplificados y compiladores modulares para facilitar la enseñanza práctica del tema \cite{navas2022modular}.

En este contexto, la visualización de los procesos internos del compilador ha emergido como una estrategia pedagógica fundamental, permitiendo a los estudiantes explorar la relación entre la teoría y la implementación mediante representaciones gráficas interactivas de autómatas y árboles sintácticos \cite{delvado2023itt}. Si bien existen herramientas educativas \cite{Stamenković2024AWE}, en el entorno de la Universidad del Valle aún se evidencia una limitada disponibilidad de recursos interactivos que permitan a los estudiantes manipular y observar de forma práctica y dinámica la evolución del ambiente conceptual de los programas.

La ausencia de recursos interactivos y accesibles localmente representa un desafío para la comprensión profunda de los contenidos de la asignatura \textit{Fundamentos de Interpretación y Compilación de Lenguajes de Programación}. La integración de enfoques visuales e interactivos en la enseñanza de compiladores ha demostrado ser una estrategia efectiva en entornos educativos \cite{Stamenković2024AWE}, lo que subraya la necesidad de explorar alternativas que promuevan una participación activa del estudiante y enriquezcan la formación práctica en este campo. 
